\section{Introduction: a quiet question in the committee room}
\label{sec:intro}

This review is written for a reader who shapes medical and artificial
intelligence law but who has never driven a surgical robot or written a line of
verification code. It therefore defines its terms as it goes, and it asks one
quiet question in eight different ways: how did we ever ask human beings, alone,
to carry a task this large without help?

A few definitions make the rest of the story readable. \textit{Physical AI} is
artificial intelligence that acts in the physical world through a machine, such
as a surgical humanoid in an oncology trial, rather than only returning text on a
screen. \textit{Verification before generation} is the principle at the center of
H. R. 9510: a machine may propose an action, but the proposal is checked against
fixed, published standards before anything touches a patient. \textit{Claude
Code} is the software agent that writes the candidate procedure; \textit{Codex}
is a second, independent agent that reviews it; and \textit{VVUQ} (verification,
validation, and uncertainty quantification) is the ten-gate examination the
proposal must pass, resolving to ACCEPT, ESCALATE to a qualified human, or BLOCK
before execution.

[DRAFTING INSTRUCTIONS: in the full-narrative, expand these definitions into a
full page using the exact sources: (a) the bill itself,
\texttt{review/references/author\_works.bib} entry \texttt{Kawchak2026HR9510Bill}
(DOI 10.5281/zenodo.20619762); (b) the verification evidence in
\texttt{cancer-automated/papers/VVUQ-02} (the ten gates, 172 of 172 automated
tests, the ACCEPT / ESCALATE / BLOCK decisions); (c) the colored Mermaid figures
in \texttt{review/mermaid/diagrams/01-verification-before-generation.md} and
\texttt{02-eight-emotional-pillars.md}. Define each lay term the first time it
appears and state the three guiding questions verbatim.]

\begin{narrativefig}
\node[nfone] (a) at (0,0) {Human specification};
\node[nftwo] (b) at (3.7,0) {Claude Code\\generate};
\node[nftwo] (c) at (7.4,0) {Codex\\review};
\node[nfdec] (d) at (10.9,0) {VVUQ\\gates};
\draw[nfedge] (a)--(b); \draw[nfedge] (b)--(c); \draw[nfedge] (c)--(d);
\end{narrativefig}
\vspace{-0.2cm}
\figcaption{Figure 1. Verification before generation, in brief. The full figure
(Mermaid 01, reproduced as colored TikZ) adds the ACCEPT, ESCALATE, and BLOCK
outcomes in the full-narrative stage.}

The three guiding questions return throughout: How did we ever utilize humans
alone for this challenging task? Why did we ever fully trust the prior human
clinical trial system if there were so many opportunities for compounding human
error? And why would the human peer-review process continue to be used for
increasingly complex and voluminous AI outputs? The eight sections that follow
answer them, in turn, through compassion, fear of preventable harm, moral
outrage, hope, responsibility, protection, trust, and urgency.
